\chapter{CONCLUSION}
\section{Introduction}This chapter walks through how LexiFlow came together, what it produced, and how close the project landed to its original goals. I'll focus on the findings, the wins worth talking about, and a few spots that could use work down the line. Basically, it's a look at how well LexiFlow actually answered the worries of teachers and learners, plus where it could grow and scale from here. The project pulled off something solid, a language learning system that genuinely met what users were hoping for. And the architecture and tech choices made along the way? They're a strong foundation for whatever comes next. With steady improvement and ongoing work, LexiFlow could end up being one of the better language learning tools out there, for universities and beyond, reliable and easy to use.
\section{Achievement of Project Objectives}The LexiFlow project started out with four main goals, which I laid out back in Section~\ref{sec:obj} of this paper. Each one was meant to make language learning feel better, partly through a friendly interface and partly through a solid audio processing pipeline. \begin{enumerate}[label= (\alph*)] \item \textbf{Elicit full Requirements from Stakeholders:} The first goal was pulling detailed requirements from stakeholders. That happened through a round of interviews with them. Those conversations gave me a real sense of the language learning needs and the headaches people run into across the UTM community. The requirements I gathered turned into the project's functional and non-functional specifications. The Software Requirements Specification (SRS) in Appendix~\ref{app:srs} documents all of it. So this one's done, the requirements got collected and analysed properly. \item \textbf{Design a Web-Based Platform:} The second goal was about designing a web-based platform with the right architecture behind it. LexiFlow was built around a Pipe and Filter architecture model. Using Vue on the frontend let me put together an intuitive, responsive interface, while Python and Node.js handled the backend services. The Software Design Document (SDD) in Appendix~\ref{app:sdd} holds the detailed design and documentation for the whole system. So, another one achieved. \item \textbf{Develop the Full Platform:} The third goal was building the language learning platform itself. That meant putting the audio processing pipeline together, which covers voice activity detection, automatic speech recognition, translation, and pinyin generation. Every feature came straight from the requirements I'd elicited and the design specs. \item \textbf{Evaluate and Test the Functionality:} The last goal was evaluating and testing how well LexiFlow actually improves the language learning experience. That part's going to happen through heavy testing, unit tests, integration tests, and user acceptance tests, in the next phase, which is PSM2. The Software Testing Documentation already holds the preliminary test cases and test approaches, attached in the appendices (Appendix~\ref{app:std}). This objective gets marked complete once the project's fully developed and the testing is done.
\section{Service Enhancement and Practical Value}So LexiFlow does more than just hit the goals it set out to. It scores on two of the three dimensions people usually use to judge whether a system in fact matters in practice: it makes the service better, and it makes life better for the people using it. Commercial value? Not yet, nobody's tried a pricing model, payment setup, or any real market testing. Here's how it improves quality of service. The stakeholder interview (Appendix~\ref{app:stakeholder-interview-questions}) flagged something painful: every simple three-minute clip of authentic Mandarin content took 10 to 30 minutes of hand work to transcribe, translate, and annotate with pinyin. LexiFlow turns that into a pipeline a learner or teacher can fire off whenever they desire and a crowd-correction feature means quality keeps climbing without anyone redoing that grind. That's a real step up from what the stakeholder described, where authentic content barely got used seeing as prepping it cost so much. And the quality-of-life side got confirmed independently during UAT (Appendix~\ref{app:uat-report}). One struggling learner figured the platform probably would've nudged their actual exam result higher. A professional in a Mandarin-speaking workplace called it something they'd reach for daily to cut down on real miscommunication at work. Neither classroom drills, generic apps, nor textbooks had touched those problems for them. These aren't vague, assumed perks, they're improvements users reported themselves, in their own learning and their everyday confidence talking to people.
\section{Suggestions for Future Improvement}So, building on what we learned from automated testing and UAT, here's what I'd recommend fixing in future versions of LexiFlow. First up: the two confirmed defects from test execution. There's the onboarding tour you can't dismiss (LF-BUG-01) and the missing SPA fallback route, which means handling directly to a URL throws back a raw 404 (LF-BUG-02). Both need fixing before any production rollout, since one blocks core navigation for first-time users and the other degrades it for returning ones. Second, there's the accuracy gap our native-speaker UAT participant flagged between LexiFlow's self-hosted faster-whisper and LibreTranslate pipeline and YouTube's commercial-grade auto-captioning. That gap should keep closing. The crowd-reporting and AI-assisted correction feature already built in Chapter 5 is set up nicely to do exactly that over time, and honestly its effectiveness deserves to be measured directly in a later evaluation round. Third, and this is the big one. A classroom-management module. Something that lets an instructor spin up a class, assign specific videos or quizzes as coursework, and check a lecturer-facing progress dashboard across all their students. The instructor stakeholder named this as their key unmet need during UAT, and since it was never part of the original functional requirements, it's the obvious pick for the next phase. Past those evidence-backed priorities, swapping the large multilingual model for a lighter Mandarin-only local one would cut processing time and server VRAM, which trims operating costs. Moving to leaner runtimes like ONNX or TensorRT would shave off even more time and resources. Taken together, these tweaks would craft LexiFlow more accurate, more usable for institutions, and cheaper to run.
